\subsection*{Exemples de funcionament del mòdul Parser}
\label{sec:parser-code-examples}

En aquest apartat es presenten diversos exemples utilitzats per comprovar el
funcionament actual del mòdul \texttt{Parser}. Cada exemple inclou el codi font
d'entrada i una representació resumida dels elements identificats. Per facilitar-ne
la lectura, s'han omès alguns nodes interns o repetitius de les sortides, com ara
els identificadors i determinats nodes de tipus, sempre que la seva inclusió no
sigui necessària per entendre l'estructura reconeguda. Per tant, les sortides que
es mostren no són una transcripció literal, sinó una síntesi de l'arbre generat.

\subsubsection*{Declaracions de tipus}
\label{sec:parser-type-declarations}

Aquest exemple permet comprovar el reconeixement de diverses declaracions de tipus
del llenguatge C\#: una classe, una estructura, una enumeració i un record. La
classe \texttt{Vehicle} també conté propietats, un camp, un constructor i un mètode,
mentre que el mètode \texttt{Main} inclou la creació d'un objecte i una invocació.

\begin{lstlisting}[style=code,
    caption={Fragment del fitxer TypeDeclarationsExample.cs},
    label={lst:type-declarations-source}]
using System;

namespace CodeSamples;

public class Vehicle
{
    public string Brand { get; set; }
    public string Model { get; set; }
    public int Year { get; set; }
    private string _privateField = "This is a private field";

    public Vehicle(string brand, string model, int year)
    {
        Brand = brand;
        Model = model;
        Year = year;
    }

    public void DisplayInfo()
    {
        Console.WriteLine($"Vehicle: {Year} {Brand} {Model}");
    }

    public static void Main()
    {
        Vehicle vehicle = new Vehicle("Volvo", "XC60", 2024);
        vehicle.DisplayInfo();
    }
}

public struct VehicleDimensions
{
    public int Length { get; set; }
    public int Width { get; set; }
}

public enum VehicleKind
{
    Car = 1,
    Truck = 2,
    Motorcycle = 3
}

public record VehicleSummary(string Brand, string Model);
\end{lstlisting}

El codi font complet d'aquest exemple es pot consultar
a \href{https://github.com/PauDiazM/static-code-analysis-tfg/blob/main/code_samples/examples/TypeDeclarationsExample.cs}{TypeDeclarationsExample.cs}.

La sortida següent resumeix els elements principals identificats pel parser:

\begin{lstlisting}[style=output,
    caption={Sortida resumida del mòdul Parser per a les declaracions de tipus},
    label={lst:type-declarations-output}]
CompilationUnit
  Import: System
  Namespace: CodeSamples
    Class: Vehicle
      Property
      Property
      Property
      Field
        VariableDeclaration
          VariableDeclarator
            EqualsValue
              Literal: This is a private field
      Callable: Vehicle(public)
        ParameterList
        Block
          Assignment
          Assignment
          Assignment
      Callable: DisplayInfo(public): void
        Block
          Invocation
            MemberAccess: Console.WriteLine
            ArgumentList
              Argument
                InterpolatedString
      Callable: Main(public static): void
        Block
          LocalDeclaration
            VariableDeclaration
              VariableDeclarator
                ObjectCreation: Vehicle
          Invocation
            MemberAccess: vehicle.DisplayInfo
    Struct: VehicleDimensions
      Property
      Property
    Enum: VehicleKind
      EnumMember: Car
        Literal: 1
      EnumMember: Truck
        Literal: 2
      EnumMember: Motorcycle
        Literal: 3
    Record: VehicleSummary
      ParameterList
        Parameter: Brand: string
        Parameter: Model: string
\end{lstlisting}

En aquest cas, el parser identifica les quatre declaracions de tipus i diversos
membres de la classe. També reconeix les assignacions del constructor, la invocació
a \texttt{Console.WriteLine}, la creació d'una instància de \texttt{Vehicle} i els
membres de l'enumeració.

\subsubsection*{Herència i classes abstractes}
\label{sec:parser-inheritance}

Aquest exemple comprova el reconeixement d'una classe abstracta i d'una classe que
n'hereta. També permet observar la representació dels mètodes abstractes i de les
implementacions que utilitzen el modificador \texttt{override}.

\begin{lstlisting}[style=code,
    caption={Codi font de l'exemple d'herència i classes abstractes},
    label={lst:inheritance-source}]
namespace CodeSamples;

public abstract class Animal
{
    public string Name { get; set; }
    public int ExtremityCount { get; set; }
    public abstract void MakeSound();
}

public class Dog : Animal
{
    public override void MakeSound()
    {
        Console.WriteLine("Woof");
    }
}
\end{lstlisting}

El codi font complet d'aquest exemple es pot consultar
a \href{https://github.com/PauDiazM/static-code-analysis-tfg/blob/main/code_samples/examples/AbstractClassExample.cs}{AbstractClassExample.cs}.

\begin{lstlisting}[style=output,
    caption={Sortida resumida del mòdul Parser per a l'exemple d'herència},
    label={lst:inheritance-output}]
CompilationUnit
  Namespace: CodeSamples
    Class: Animal
      Property
      Property
      Callable: MakeSound(public abstract): void
        ParameterList
    Class: Dog
      InheritanceList
        InheritedType: Animal
      Callable: MakeSound(public override): void
        ParameterList
        Block
          ExpressionStatement
            Invocation
              MemberAccess: Console.WriteLine
              ArgumentList
                Argument
                  Literal: Woof
\end{lstlisting}

La sortida mostra que la relació d'herència es representa mitjançant els nodes
\texttt{InheritanceList} i \texttt{InheritedType}. Així mateix, els mètodes de
les dues classes es representen com a elements \texttt{Callable}, conservant els
modificadors \texttt{abstract} i \texttt{override}.

\subsubsection*{Estructures de control condicionals}
\label{sec:parser-conditionals}

Aquest exemple conté una estructura \texttt{if}, una branca \texttt{else if} i una
branca \texttt{else}. També inclou declaracions de variables, expressions binàries
i assignacions dins dels diferents blocs condicionals.

\begin{lstlisting}[style=code,
    caption={Codi font de l'exemple d'estructures condicionals},
    label={lst:conditional-source}]
public class ConditionalExample
{
    public static void Main()
    {
        int score = 75;
        string result = "Unknown";

        if (score >= 90)
        {
            result = "Excellent";
        }
        else if (score >= 60)
        {
            result = "Passed";
        }
        else
        {
            result = "Try again";
        }

        Console.WriteLine(result);
    }
}
\end{lstlisting}

El codi font complet d'aquest exemple es pot consultar
a \href{https://github.com/PauDiazM/static-code-analysis-tfg/blob/main/code_samples/examples/ConditionalExample.cs}{ConditionalExample.cs}.


\begin{lstlisting}[style=output,
    caption={Sortida resumida del mòdul Parser per a les estructures condicionals},
    label={lst:conditional-output}]
CompilationUnit
  Class: ConditionalExample
    Callable: Main(public static): void
      Block
        LocalDeclaration
          VariableDeclaration
            VariableDeclarator
              EqualsValue
                Literal: 75
        LocalDeclaration
          VariableDeclaration
            VariableDeclarator
              EqualsValue
                Literal: Unknown
        If
          BinaryExpression: >=
            Identifier: score
            Literal: 90
          Block
            ExpressionStatement
              Assignment
                Identifier: result
                Literal: Excellent
          Else
            If
              BinaryExpression: >=
                Identifier: score
                Literal: 60
              Block
                ExpressionStatement
                  Assignment
                    Identifier: result
                    Literal: Passed
              Else
                Block
                  ExpressionStatement
                    Assignment
                      Identifier: result
                      Literal: Try again
        ExpressionStatement
          Invocation
            MemberAccess: Console.WriteLine
            ArgumentList
              Argument
                Identifier: result
\end{lstlisting}

Cal destacar que la construcció \texttt{else if} es representa com un node
\texttt{If} niat dins del node \texttt{Else} corresponent. Aquesta representació
permet mantenir l'estructura jeràrquica de les branques condicionals.

\subsubsection*{Invocacions i accés a membres}
\label{sec:parser-invocations}

Aquest exemple, de dimensions reduïdes, permet observar com es representa una
invocació de mètode i l'accés al membre corresponent. També mostra el tractament
dels arguments i dels literals de tipus cadena.

\begin{lstlisting}[style=code,
    caption={Codi font de l'exemple d'invocació},
    label={lst:invocation-source}]
class HelloWorld
{
    static void Main()
    {
        Console.WriteLine("Hello World!");
    }
}
\end{lstlisting}

El codi font complet d'aquest exemple es pot consultar
a \href{https://github.com/PauDiazM/static-code-analysis-tfg/blob/main/code_samples/examples/HelloWorld.cs}{HelloWorld.cs}.

\begin{lstlisting}[style=output,
    caption={Sortida resumida del mòdul Parser per a una invocació},
    label={lst:invocation-output}]
CompilationUnit
  Class: HelloWorld
    Callable: Main(static): void
      ParameterList
      Block
        ExpressionStatement
          Invocation
            MemberAccess: Console.WriteLine
            ArgumentList
              Argument
                Literal: Hello World!
\end{lstlisting}

En aquest cas, la invocació es representa mitjançant un node \texttt{Invocation},
que conté l'accés al membre \texttt{Console.WriteLine} i la llista d'arguments.

\subsubsection*{Declaracions d'interfícies}
\label{sec:parser-interfaces}

Aquest exemple comprova el reconeixement de les interfícies i dels mètodes que s'hi
declaren sense proporcionar-ne una implementació.

\begin{lstlisting}[style=code,
    caption={Codi font de l'exemple d'interfície},
    label={lst:interface-source}]
public interface Game
{
    void Start();
    void Stop();
    void Pause();
}
\end{lstlisting}

El codi font complet d'aquest exemple es pot consultar
a \href{https://github.com/PauDiazM/static-code-analysis-tfg/blob/main/code_samples/examples/InterfaceExample.cs}{InterfaceExample.cs}.

\begin{lstlisting}[style=output,
    caption={Sortida resumida del mòdul Parser per a una interfície},
    label={lst:interface-output}]
CompilationUnit
  Interface: Game
    Callable: Start(): void
      ParameterList
    Callable: Stop(): void
      ParameterList
    Callable: Pause(): void
      ParameterList
\end{lstlisting}

Els mètodes declarats a la interfície es representen amb el mateix tipus genèric
\texttt{Callable} que s'utilitza per als mètodes de les classes. D'aquesta manera,
es manté una representació comuna per a les declaracions de mètodes, independentment
del tipus que les conté.

\subsubsection*{Estructures iteratives i arrays}
\label{sec:parser-loops-arrays}

L'últim exemple inclou les quatre estructures iteratives principals de C\#:
\texttt{for}, \texttt{foreach}, \texttt{while} i \texttt{do-while}. També permet
comprovar el reconeixement de tipus array, inicialitzadors, accessos a elements i
expressions d'increment.

\begin{lstlisting}[style=code,
    caption={Codi font de l'exemple d'estructures iteratives i arrays},
    label={lst:loops-source}]
public class LoopTests
{
    public static void Main(string[] args)
    {
        // Common C# loop
        for (int i = 0; i < 5; i++)
        {
            Console.WriteLine($"Iteration {i + 1}");
        }

        // Foreach loop
        string[] fruits = { "Apple", "Banana", "Cherry" };
        Console.WriteLine(fruits[0]);
        foreach (string fruit in fruits)
        {
            Console.WriteLine(fruit);
        }

        // While loop
        int count = 0;
        while (count < 3)
        {
            count++;
        }

        // Do-while loop
        int doCount = 0;
        do
        {
            doCount++;
        } while (doCount < 2);
    }
}
\end{lstlisting}

El codi font complet d'aquest exemple es pot consultar
a \href{https://github.com/PauDiazM/static-code-analysis-tfg/blob/main/code_samples/examples/LoopExample.cs}{LoopExample.cs}.

\begin{lstlisting}[style=output,
    caption={Sortida resumida del mòdul Parser per a les estructures iteratives},
    label={lst:loops-output}]
CompilationUnit
  Class: LoopTests
    Callable: Main(public static): void
      ParameterList
        Parameter: args: string[]
          ArrayType
            ArrayDimensionList
              ArrayDimension
      Block
        For
          VariableDeclaration
            VariableDeclarator
              EqualsValue
                Literal: 0
          BinaryExpression: <
            Identifier: i
            Literal: 5
          PostfixUnaryExpression: ++
            Identifier: i
          Block
            ExpressionStatement
              Invocation
                MemberAccess: Console.WriteLine
                ArgumentList
                  Argument
                    InterpolatedString
                      Interpolation
                        BinaryExpression: +
        LocalDeclaration
          VariableDeclaration
            ArrayType
              ArrayDimensionList
                ArrayDimension
            VariableDeclarator
              EqualsValue
                Initializer
                  Literal: Apple
                  Literal: Banana
                  Literal: Cherry
        ExpressionStatement
          Invocation
            MemberAccess: Console.WriteLine
            ArgumentList
              Argument
                ElementAccess: fruits
                  ArgumentList
                    Argument
                      Literal: 0
        ForEach
          PredefinedType
          Identifier: fruits
          Block
            ExpressionStatement
              Invocation
                MemberAccess: Console.WriteLine
                ArgumentList
                  Argument
                    Identifier: fruit
        LocalDeclaration
          VariableDeclaration
            VariableDeclarator
              EqualsValue
                Literal: 0
        While
          BinaryExpression: <
            Identifier: count
            Literal: 3
          Block
            ExpressionStatement
              PostfixUnaryExpression: ++
                Identifier: count
        LocalDeclaration
          VariableDeclaration
            VariableDeclarator
              EqualsValue
                Literal: 0
        DoWhile
          Block
            ExpressionStatement
              PostfixUnaryExpression: ++
                Identifier: doCount
          BinaryExpression: <
            Identifier: doCount
            Literal: 2
\end{lstlisting}

La sortida mostra que cada estructura iterativa disposa d'un node específic. També
s'hi poden observar els nodes utilitzats per representar les dimensions dels arrays,
la seva inicialització i l'accés a un dels seus elements. En el cas del bucle
\texttt{for}, la inicialització, la condició i l'actualització apareixen com a
components diferenciats del node corresponent.